% ID: SEL-D0023-OA-E03
% ESTADO: cotejada. Segunda pasada visual completa y borrador compilado revisados.
% Fuente: Selectividad/matematicas_aplicadas/reserva_b_EXAMEN 2 -ANDALUCÍA.pdf
% SHA-256: 8cdb86259405652622f2ea129bffbdb239076ac612fd7fcebb86b9d5bdc451ec
% Página: 1; opción A; ejercicio original 3.

En un centro de estudios que tiene 250 estudiantes, hay 50 que tienen problemas visuales y 20 que tienen problemas auditivos. Los sucesos ``tener problemas visuales'' y ``tener problemas auditivos'' son independientes.

Se elige un estudiante al azar, calcule las probabilidades de los sucesos siguientes:
\begin{enumerate}[label=\alph*)]
\item \textbf{(0.75 puntos)} Tener problemas visuales y auditivos.
\item \textbf{(0.75 puntos)} No tener problemas visuales ni auditivos.
\item \textbf{(1 punto)} Tener algún problema auditivo si no tiene problemas visuales.
\end{enumerate}
